\section{Outlook}\label{sec:outlook}

\subsection{What is settled}\label{sec:settled}

Problem~\ref{prob:main} is answered: $d=2\omega$ is a parameter. The dictionary
names one object, a rank-$(d{+}1)$ lattice (Section~\ref{sec:rank}); the object
exists at every integer and nothing strictly in between
(Section~\ref{sec:pointcount}); the transfer continues past the family's range
with its decomposition and its positivity intact
(Proposition~\ref{prop:cmuncapped}) while its criterion does not
(Proposition~\ref{prop:vacuity}); and the transfer's positivity continues for
reasons that have nothing to do with a lattice
(Proposition~\ref{prop:position}), so the lattice's failure leaves no trace in
it. There is no remaining sense in which $d$ is a dimension away from the
integers, and at fractional $d$ there is a scattering matrix without a space.

Three questions that earlier stages of the investigation posed are closed
negatively, and the closures are worth as much as the openings.

\begin{itemize}
\item \emph{Does the criterion continue to the half-integer points beyond the
 range, where the object is a genuine lattice?} It continues, and the
 continuation is vacuous (Proposition~\ref{prop:vacuity}). Testing contraction at
 $\omega=1,\frac32,2$ cannot fail, so it calibrates the instrument and learns
 nothing about $\zeta$.
\item \emph{Is the exchange of complete monotonicity for innerness at $m=2$ the
 same fact as the point count's failure?} No (Corollary~\ref{cor:noequiv}); the
 shared interval is the first gap (Remark~\ref{rem:forcedagreement}).
\item \emph{Is there a factorization of the transfer that exhibits which factor
 loses positivity at fractional $d$?} No (Corollary~\ref{cor:nofactorization}),
 because no factor loses positivity anywhere.
\end{itemize}

\subsection{What is open}\label{sec:open}

\paragraph{The complex-hyperbolic line.} In the true half-distance the
archimedean measure is the radial measure of complex hyperbolic space of complex
dimension $\omega$ and real dimension $2\omega$, with Jacobi parameters
$(\alpha,\beta)=(\omega-1,0)$ and $\rho=\omega>0$ \cite{WhichDimension} --- a
Heckman--Opdam parameter, where a transform theory exists at arbitrary
multiplicity \cite{HeckmanSchlichtkrull}. Proposition~\ref{prop:position} now
explains why that half of the transfer continues at all: its positivity is a Beta
density with $\omega$ in the exponent. The question is whether the Heckman--Opdam
theory supplies anything beyond a name for that --- an inversion, a Plancherel
measure, a completeness statement that the Beta density does not already give.
The honest prior is that it supplies a name.

\paragraph{Two conjectures about sums of squares.} That $r_m(n)\ge0$ for every
$n$ and every real $m\ge4$; and that the negative coefficients on $m\in(3,4)$
arise from competition between terms of both signs rather than from forcing
(Remark~\ref{rem:lagrange}). Neither is a question about $\zeta$, and neither is
used here.

\paragraph{A convention to retire.} Remark~\ref{rem:xiwarning}: three research
notes of this investigation write $\xi$ for what this manuscript and its parent
call $\Lambda$, and the two differ by exactly $R_\omega$.

\subsection{Status of every statement}\label{sec:status}

\begin{itemize}
\item \textbf{Written proof, unconditional:}
 Propositions~\ref{prop:epstein}, \ref{prop:integerpoints},
 \ref{prop:twoblaschke}, \ref{prop:cmuncapped}, \ref{prop:vacuity},
 \ref{prop:signlaw}, \ref{prop:free}, \ref{prop:Nk}, \ref{prop:forced},
 \ref{prop:onethreshold}, \ref{prop:position}; and
 Corollaries~\ref{cor:content}, \ref{cor:nolattice}, \ref{cor:noequiv},
 \ref{cor:nofactorization}. None is conditional on the Riemann hypothesis and
 none mentions a zero except to bound its real part by the Euler product.
\item \textbf{Written proof citing a classical theorem:}
 Proposition~\ref{prop:Nk}, whose last clause is Lagrange's four-square theorem
 \cite{HardyWright}, used and not reproved; and Remark~\ref{rem:boundary}, which
 uses the classical zero-free region \cite{Titchmarsh} at the single point
 $\omega=\frac12$.
\item \textbf{Inherited and used as stated:}
 Propositions~\ref{prop:unimodular}, \ref{prop:dichotomy},
 \ref{prop:firstorder}, \ref{prop:pieces} and~\ref{prop:factor}, all from
 \cite{Markov}; and Proposition~\ref{prop:interp} from \cite{RankNote}.
\item \textbf{Registered computation:} the six programmes of
 Appendix~\ref{app:numerics}, $1522$ finite test cases, each with a preserved
 record and each verified deterministic across repeated runs. Counts refer to
 test cases, not to independent theorems.
\item \textbf{Calibration, not evidence:} every number in
 Table~\ref{tab:calibration}, for the reason in Remark~\ref{rem:direction}.
\item \textbf{Reading:} the group-theoretic sentence of
 Corollary~\ref{cor:oneobject}; and Remark~\ref{rem:reading}.
\item \textbf{Conjecture, used nowhere:} the two statements of
 Remark~\ref{rem:lagrange}.
\item \textbf{Not claimed:} that any object exists at fractional $d$; that the
 transfer is a contraction at any $\omega<\frac12$; anything about the location
 of the zeros; any positivity of the Weil form.
\end{itemize}
